\chapter{JWT Authentication}

Once a password is verified, the server needs to remember the browser is logged in on every later request -- without a database lookup each time. A \textbf{JSON Web Token (JWT)} packs identity into a signed, self-contained string the server can verify without storing anything.

\section{Anatomy of a JWT}

A JWT is three base64url segments separated by dots: \texttt{header.payload.signature}.

\begin{center}
\begin{tabularx}{\textwidth}{lX}
\toprule
\textbf{Segment} & \textbf{Contents} \\
\midrule
Header & Signing algorithm (e.g. \texttt{HS256}) and token type. \\
Payload & Claims, e.g. \texttt{\{ id: userId, iat, exp \}} -- readable by anyone, not encrypted. \\
Signature & \texttt{HMACSHA256(header + payload, JWT\_SECRET)} -- proves origin and integrity. \\
\bottomrule
\end{tabularx}
\end{center}

\begin{warnbox}[WARNING]
The payload is base64-encoded, not encrypted -- anyone can decode it. Never put passwords or sensitive data in a JWT payload, only identifiers.
\end{warnbox}

\subsection{Expiration and Signing}

Every token needs an \texttt{exp} claim so a leaked token becomes useless after a fixed window. The signature uses \texttt{process.env.JWT\_SECRET}, known only to your server -- treat it like a password.

\section{The Authentication Flow}

\begin{diagrambox}[JWT Authentication Flow]
\begin{verbatim}
 1. Browser -> POST /api/auth/login (email, password)
 2. Express -> verify credentials (bcrypt.compare)
 3. Express -> generate JWT (jwt.sign, JWT_SECRET)
 4. Express -> send JWT in HttpOnly cookie
 5. Browser -> stores cookie, sends it on every request
 6. Browser -> GET /api/tasks  (cookie attached automatically)
 7. Middleware -> jwt.verify -> req.user = user document
 8. Controller -> runs, trusting req.user
\end{verbatim}
\end{diagrambox}

\section{Generating a Token}

\begin{lstlisting}[language=JavaScript]
// utils/generateToken.js
import jwt from "jsonwebtoken";

export const generateToken = (userId) => {
  return jwt.sign(
    { id: userId },
    process.env.JWT_SECRET,
    { expiresIn: "7d" }
  );
};
\end{lstlisting}

Used right after credentials are verified in the login controller:

\begin{lstlisting}[language=JavaScript]
// controllers/authController.js
import { generateToken } from "../utils/generateToken.js";

export const loginUser = async (req, res, next) => {
  try {
    const { email, password } = req.body;
    const user = await User.findOne({ email });
    if (!user || !(await bcrypt.compare(password, user.password))) {
      return res.status(401).json({ success: false, message: "Invalid credentials" });
    }

    const token = generateToken(user._id);

    res.cookie("token", token, {
      httpOnly: true,
      maxAge: 7 * 24 * 60 * 60 * 1000,
    });

    res.status(200).json({
      success: true,
      data: { id: user._id, name: user.name, role: user.role },
    });
  } catch (err) {
    next(err);
  }
};
\end{lstlisting}

\begin{notebox}[NOTE]
Full cookie config (\texttt{secure}, \texttt{sameSite}, why \texttt{httpOnly} matters) is covered next chapter.
\end{notebox}

\section{Verifying a Token}

\begin{lstlisting}[language=JavaScript]
// middleware/authMiddleware.js
import jwt from "jsonwebtoken";

export const protect = async (req, res, next) => {
  const token = req.cookies.token;
  if (!token) {
    return res.status(401).json({ success: false, message: "Not authenticated" });
  }

  try {
    const decoded = jwt.verify(token, process.env.JWT_SECRET);
    req.userId = decoded.id; // full lookup covered in Chapter 20
    next();
  } catch (err) {
    return res.status(401).json({ success: false, message: "Invalid or expired token" });
  }
};
\end{lstlisting}

\texttt{jwt.verify} checks the signature and rejects an expired \texttt{exp} claim automatically -- both throw, hence the \texttt{try/catch}.

\begin{tipbox}[Advanced: Access Tokens vs Refresh Tokens]
Larger systems split this into a short-lived access token plus a long-lived refresh token used only to silently reissue access tokens, limiting damage if one leaks at the cost of extra complexity. This book uses a single 7-day token to stay focused.
\end{tipbox}

\whatwhyhow
{A JWT is a signed token containing a user identifier and expiration, generated on login and verified on every protected request.}
{It lets the server confirm "who is this request from" without a database lookup on every request or server-side session storage, while still expiring automatically.}
{Sign with \texttt{jwt.sign(\{id\}, JWT\_SECRET, \{expiresIn\})} at login, store it in a cookie, and verify with \texttt{jwt.verify(token, JWT\_SECRET)} inside middleware before any protected controller runs.}
